% SOURCE_AUTHORITY: sga5_fr_workpass.tex, lines 9656--9692 (LF-normalized unit).
% SOURCE_SHA256: 791F4EFFC5E02832D5D77ED03518C8156D6F07E4C8238B03545DB93D883FBB28
\emph{Demonstração.} Basta verificar quando $E$ é inversível. Então, com as identificações usuais, $\varepsilon^*(E)$ é a imagem de $E$ pela canônica aplicação.
\[
\varepsilon^*:H^1(X_{\mathrm{\acute et}},\mathbb G_m)\to H^1(X(\mathbb C),\mathcal O_c^*),
\tag{3.8.6}
\]
e é uma questão de ver que o diagrama
\[
\begin{tikzcd}[column sep=large,row sep=large]
H^1(X,\mathbb G_m) \arrow[r,"\varepsilon^*"] \arrow[d,"\partial"'] &
H^1(X(\mathbb C),\mathcal O_c^*) \arrow[d,"c_1"]\\
H^2(X,\mu_X) \arrow[r,"\varepsilon^*"] &
H^2(X(\mathbb C),\mathbb Z/\nu\mathbb Z)
\end{tikzcd}
\tag{3.8.7}
\]
onde $\partial$ é o morfismo de borda proveniente da sequência exata canônica
\[
0\to \mu_X\to\mathbb G_m\xrightarrow{\nu}\mathbb G_m\to0
\tag{3.8.8}
\]
e $c_1$ a primeira classe de Chern (3.8.1), é comutativo. De fato, usando que o morfismo primeiro classe de Chern
\[
H^1(X(\mathbb C),\mathcal O_\mathbb C^*)\to H^2(X(\mathbb C),\mathbb Z)
\]
é o morfismo da borda da sequência exata canônica
\[
0\to\mathbb Z\to\mathcal O_c\xrightarrow{\exp}\mathcal O_c^*\to0,
\qquad
t\mapsto 2i\pi t,
\]
É fácil de ver que o morfismo $c_1$ de (3.8.7) não é outro que o morfismo de borda associado à sequência exata
\[
0\to\mathbb Z/\nu\mathbb Z\to\mathcal O_c^*\xrightarrow{\nu}\mathcal O_c^*\to0,
\qquad
t\mapsto e^{\frac{2i\pi t}{\nu}},
\]
que, tendo em conta (3.8.4), é o análogo de (3.8.8). A partir daí segue-se imediatamente a afirmação.
